%中文字体均设置为思源相应字体,若未安装思源宋体则使用下行命令,并酌情删除代码
%\usepackage[fontset=fandol]{ctex}
\usepackage[fontset=none]{ctex}
\usepackage{xeCJK}
\xeCJKsetup{CJKmath=true} %数学模式中可以输入中文
\setCJKmainfont{思源宋体}[ItalicFont=FandolKai]
\setCJKsansfont{思源黑体}
\setCJKmonofont{FandolFang}

%AMS全家桶,\DeclareMathOperator依赖之
\usepackage{amsmath,amssymb,amsthm,amsfonts}
\usepackage{tikz-cd} %用来画交换图
\usepackage{bm} %粗体字母(含希腊字母)

%纸张上下左右页边距
\usepackage[a4paper,left=1cm,right=1cm,top=1.5cm,bottom=1.5cm]{geometry}
%生成书签和目录上的超链接
\usepackage[colorlinks=true,linkcolor=blue,filecolor=blue,urlcolor=blue,citecolor=cyan]{hyperref}

%各种列表环境的行距
\usepackage{enumitem}
\setenumerate[1]{itemsep=0pt,partopsep=0pt,parsep=\parskip,topsep=10pt}
\setenumerate[2]{itemsep=0pt,partopsep=0pt,parsep=\parskip,topsep=0pt}
\setenumerate[3]{itemsep=0pt,partopsep=0pt,parsep=\parskip,topsep=0pt}
\setitemize[1]{itemsep=0pt,partopsep=0pt,parsep=\parskip,topsep=5pt}
\setdescription{itemsep=0pt,partopsep=0pt,parsep=\parskip,topsep=5pt}
%在且仅在读书笔记中使得enumerate环境的数字编号加粗
\ifdefined\note
\setenumerate[1]{label=\textbf{\arabic*.}}
\else\fi

\everymath{\displaystyle} %全体公式为行间形式
\setmainfont{CMU Serif} %能显示大量外文字体
% 若在param中新定义了命令\dayhruletrue,则此命令生成分割线,否则不生成
\newcommand{\dayhrule}{\ifdefined\dayhruletrue\hrulefill\else\fi}

%收敛箭头
\newcommand{\converge}[1]{\xrightarrow{#1}}
\newcommand{\convergedot}{\converge{\cdot}}
\newcommand{\convergeuni}{\converge{\mathrm{uni.}}}
%左右配对符号
\newcommand{\br}[1]{\!\left(#1\right)} %括号
\newcommand{\cbr}[1]{\left\{#1\right\}} %大括号
\newcommand{\abr}[1]{\left<#1\right>} %尖括号(内积)
\newcommand{\bbr}[1]{\left[#1\right]} %中括号
\newcommand{\abbr}[1]{\left(#1\right]} %左开右闭区间
\newcommand{\babr}[1]{\left[#1\right)} %左闭右开区间
\newcommand{\abs}[1]{\left|#1\right|} %绝对值
\newcommand{\norm}[1]{\left\|#1\right\|} %范数
\newcommand{\floor}[1]{\left\lfloor#1\right\rfloor} %下取整
\newcommand{\ceil}[1]{\left\lceil#1\right\rceil} %上取整
%左右配对符号(标准大小)
\newcommand{\brnormal}[1]{(#1)} %括号
\newcommand{\cbrnormal}[1]{\{#1\}} %大括号
\newcommand{\abrnormal}[1]{\langle#1\rangle} %尖括号(内积)
\newcommand{\bbrnormal}[1]{[#1]} %中括号
\newcommand{\absnormal}[1]{|#1|} %绝对值
\newcommand{\normnormal}[1]{\|#1\|} %范数
%常用数集简写
\newcommand{\R}{\mathbb{R}} %实数域
\newcommand{\N}{\mathbb{N}} %自然数集
\newcommand{\Z}{\mathbb{Z}} %整数集
\newcommand{\C}{\mathbb{C}} %复数域
\newcommand{\F}{\mathbb{F}} %一般数域
\newcommand{\Q}{\mathbb{Q}} %有理数域
\renewcommand{\P}{\mathbb{P}} %本产生¶(段落符号)
%代数符号
\DeclareMathOperator{\aut}{Aut} %自同构
\DeclareMathOperator{\inn}{Inn} %内自同构
\DeclareMathOperator{\inv}{Inv}
\DeclareMathOperator{\GL}{GL} %一般线性群
\DeclareMathOperator{\GF}{GF}
\DeclareMathOperator{\SL}{SL} %特殊线性群
%正体符号
\renewcommand{\i}{\mathrm{i}} %本产生无点i
\newcommand{\id}{\mathrm{id}} %恒等映射
\newcommand{\e}{\mathrm{e}} %自然常数e
\renewcommand{\d}{\mathrm{d}} %微分符号，本产生重音符号
\newcommand{\D}{\partial} %偏导符号
\newcommand{\dfunc}[2]{\frac{\d #1}{\d #2}}
\newcommand{\Dfunc}[2]{\frac{\D #1}{\D #2}}
%运算符
\DeclareMathOperator{\Arg}{Arg} %辐角
\newcommand{\Span}{\mathrm{span}} %张成空间
\DeclareMathOperator{\re}{Re} %实部
\DeclareMathOperator{\im}{Im} %像,虚部
\DeclareMathOperator{\tr}{tr} %迹
\DeclareMathOperator{\grad}{grad} %梯度
\DeclareMathOperator{\rank}{rank} %秩
\DeclareMathOperator{\charfield}{char} %域的特征
\DeclareMathOperator{\lcm}{lcm} %最小公倍数
\DeclareMathOperator{\sgn}{sgn} %符号函数
\DeclareMathOperator{\codim}{codim} %余维度
\DeclareMathOperator{\conv}{conv} %凸包
\DeclareMathOperator{\supp}{supp} %支撑
\DeclareMathOperator{\Log}{Log} %广义对数函数
\DeclareMathOperator{\card}{card} %集合的势
\DeclareMathOperator{\Res}{Res} %留数
%简写
\newcommand{\ds}{\displaystyle} %行间公式形式
\newcommand{\ve}{\varepsilon} %手写体ε
\newcommand{\rev}{^{-1}\!} %逆
\newcommand{\T}{^{\mathsf{T}}} %转置
\renewcommand{\H}{^{\mathsf{H}}} %共轭转置
\newcommand{\adj}{^\lor} %伴随
\newcommand{\PowSet}[1]{\mathcal{P}\!\left(#1\right)} %幂集
\newcommand{\textif}{\mathrm{if}\,}
\newcommand{\sth}{\mathrm{sth.}}
\DeclareMathOperator{\lhs}{LHS}
\DeclareMathOperator{\rhs}{RHS}
\newcommand{\other}{\text{Otherwise}}
\newcommand{\hint}[1]{{\small (#1)}} %提示
\newcommand{\why}{\textcolor{red}{(Why?)}}
\newcommand{\tbc}{\textcolor{red}{(To be continued...)}} %未完待续

%定理环境(随笔记形式更改)
\theoremstyle{definition}
\newtheorem{definition}{定义}
\newtheorem{remark}{注}
\ifdefined\notebook{
    \newtheorem{example}{例}
    \theoremstyle{plain}
    \newtheorem{theorem}{定理}[chapter]
}
\else{
    \newtheorem{example}{例}[section]
    \newtheorem{theorem}{定理}[section]
}
\fi
\theoremstyle{plain}
\newtheorem{lemma}[theorem]{引理}
\newtheorem{proposition}[theorem]{命题}